This paper studies how to evaluate topology-aware federated control methods for active voltage control under topology shift on the case141 benchmark. Instead of relying on mixed-context summaries, we use a deterministic, context-aligned runtime bundle that pairs each federated method with matched baseline rollouts and reports paired seed deltas for return, voltage violation, and active-power loss. In the completed three-seed main benchmark, neither \texttt{fedgrid\_topo\_proto} nor \texttt{fedgrid\_v4\_cluster\_distill} shows a reliable advantage over the baseline on the \texttt{random\_reset} benchmark: mean paired $\Delta$return is $-0.106$ and $-0.173$, respectively, and each method wins only $1/3$ seeds. A completed one-seed robustness suite shows that stress variants can behave differently, with \texttt{fedgrid\_v4\_cluster\_byzantine} reaching $+0.642$ paired $\Delta$return on \texttt{random\_reset}, but that result is not strong enough to support a broad superiority claim because it is single-seed. The current evidence therefore supports an empirical paper about evaluation discipline, context sensitivity, and failure modes rather than a blanket new-method claim. We release the runnable training, deterministic evaluation, aggregation, table, figure, and reporting pipeline together with the current suite artifacts so the conclusions remain traceable and reproducible.
